\chapter{Flags}
System flags control PCB manufacturing capabilities using the \texttt{syst.flag.} prefix:
\begin{verbatim}
syst.flag.<flag_name> <value>
\end{verbatim}

\section{Available Flags}
\begin{itemize}
\item \texttt{exec\_mode} -- Execution mode: \texttt{lenient} (default) or \texttt{strict}. In strict mode, all warnings become errors and halt execution.
\item \texttt{vias} -- Controls whether vias are allowed
\item \texttt{slots} -- Controls whether slots / PCB cuts are allowed
\item \texttt{jumpers} -- Controls whether jumper wires are allowed
\item \texttt{solder\_reinforced\_tracks} -- Controls whether solder-reinforced tracks are allowed
\item \texttt{buried\_vias} -- Controls whether buried vias are allowed
\item \texttt{approx\_eps} -- Additive tolerance for approximate equality operator \texttt{\textasciitilde} (default: 1)
\item \texttt{approx\_eps\_mult} -- Multiplicative tolerance for approximate equality operator \texttt{\textasciitilde*} (default: 1)
\end{itemize}


\part{Toolchain}
\chapter{Compiler and Execution Pipeline}
The EDADSL toolchain compiles source code into PCB and schematic files.

\section{Compilation Phases}
\begin{enumerate}
\item Lexical Analysis -- Tokenization of source code
\item Parsing -- Syntax analysis and AST construction
\item Semantic Analysis -- Type checking and validation
\item Execution -- Program execution with electrical system construction
\item Constraint Solving -- Physical layout optimization
\item Code Generation -- Output file generation
\end{enumerate}

\section{Entry Point}
Every program must begin with the \texttt{start} keyword. The program halts when \texttt{halt} is reached or an unrecoverable error occurs.


\chapter{Code Generation}

\section{PCB Generation}
The \texttt{make.pcb} command generates a PCB file based on the circuit and physical constraints:
\begin{verbatim}
make.pcb
\end{verbatim}
This command must be placed after all circuit definitions and physical constraints.

\section{Schematic Generation}
The \texttt{make.schematic} command generates a schematic file:
\begin{verbatim}
make.schematic
\end{verbatim}

\section{Export}
The \texttt{halt} command terminates execution and exports all created files:
\begin{verbatim}
halt
\end{verbatim}


\chapter{Standard Library}
% TODO: Document the standard library contents
% Uncertain about:
% - Complete list of available component models
% - Footprint library organization
% - Material library contents


\chapter{External Library}
% TODO: Document external library integration
% Uncertain about:
% - How external libraries are imported
% - Supported library formats
% - KiCad library integration details


\chapter{Debugging and Logging}

\section{Console Output}
The \texttt{echo} and \texttt{f.print()} functions output to the console:
\begin{verbatim}
echo $variable
f.print("Message: ", $variable)
\end{verbatim}

\section{Log File}
The \texttt{writetolog} keyword appends a message to the log file:
\begin{verbatim}
writetolog "Checkpoint reached"
\end{verbatim}
The log file is exported when \texttt{halt} is executed.




\part{Reference}
\chapter{Exceptions}

\section{Type Errors}
Raised when an operation is applied to a value of the wrong type.

\section{Value Errors}
Raised when a function receives an argument outside its domain (e.g., division by zero, invalid list index). Out-of-bounds list or set indexing (e.g., accessing index 10 on a 3-element list) raises a Value Error and halts execution.

\section{Reference Errors}
Raised when referencing an undefined variable, part, or function.

\section{Constraint Errors}
Raised when physical constraints cannot be resolved (see Section~\ref{sec:solver-model}).


\chapter{Error and Warning Model}

\section{Errors}
Errors halt program execution. The error message includes the error type and location.

\section{Warnings}
Warnings are printed for constraint conflicts between different priorities. The program continues execution after warnings.


\chapter{Future Improvements}
% TODO: Document planned language features and improvements
% Uncertain about:
% - What features are planned for future versions
% - Deprecation timeline for any features


\chapter{Implementations}

\section{Reference Implementation}
% TODO: Document the reference implementation details
% Uncertain about:
% - Programming language of the implementation
% - Current version number
% - Supported platforms


\chapter{References}

\section{Language Specification}
This document serves as the authoritative specification for the EDADSL language.

\section{Related Standards}
% TODO: Reference relevant EDA standards (e.g., KiCad file formats, IPC standards)




